% Hand-edited. Sources and claim boundaries are recorded in positioning.md.
\section{Related Work}
\label{sec:related-work}

\paragraph{Post-hoc activation sparsification.}
CATS applies magnitude thresholding within gated FFNs and studies both
training-free use and subsequent fine-tuning~\citep{lee2024cats}.
TEAL applies magnitude thresholding across attention and FFN projections,
studies input versus output sparsity, and uses block-wise greedy allocation to
accommodate heterogeneous sensitivity~\citep{liu2024teal}. Thus, both the
placement of sparsity and its allocation across projections have precedents.
Our question concerns how those choices interact with activation pressure
during pretraining. Post-hoc thresholding provides a reference for the
trade-off available without changing the learned weights. Uniform
TEAL-style thresholding is a restricted reference and does not reproduce
TEAL's optimized allocation procedure.

\paragraph{Learning sparse representations.}
ReLU Strikes Back studies restoring sparse activations in language
models~\citep{mirzadeh2023relu}. ProSparse adds progressive $\ell_1$ pressure
and threshold shifting, including ablations showing that the regularization
schedule affects the quality obtained at comparable FFN
sparsity~\citep{song2024prosparse}. Sparsing Law further studies the dependence
of activation sparsity on model size, training data, and regularization
strength~\citep{luo2025sparsing}.
\citet{cetin2026sparser} combine FFN activation regularization with sparse
execution for training and inference. These works establish that sparsity
depends on the training procedure as well as the activation function. Their
FFN activation fractions, however, measure a different quantity from a product
count that includes attention and the dense output head. We build on this
literature by separating pressure from the gates and sites through which it
affects the complete forward pass.

\paragraph{Sparsity across attention and FFNs.}
Q-Sparse trains with top-$k$ selection on inputs to attention and FFN linear
projections, using a straight-through estimator and squared ReLU; it studies
training from scratch, continued training, and fine-tuning, and reports overall
sparsity and activated-parameter budgets~\citep{wang2024qsparse}. It therefore
already establishes training-time sparsification beyond FFNs. Its projection
input masks should be distinguished from masks on the query, key, and value
operands inside attention. Spark Transformer combines sparse FFNs with
top-$k$ attention~\citep{you2025spark}; selecting attention positions is also
different from zeroing feature coordinates in those operands. Our focus is
the marginal effect of pressure and gate placement under a shared training
protocol. The contribution sought is an explanation of that interaction,
including its dependence on the operation being sparsified.

\paragraph{Task and regularization gradients.}
PCGrad projects away conflicting gradient components in multi-task
learning~\citep{yu2020pcgrad}. Bloop treats an auxiliary objective as secondary
to the training objective and uses a moving average to improve gradient
surgery~\citep{hsieh2024bloop}. These methods motivate examining optimization
interference when activation pressure is added. Our orthogonal $\ell_1$
variant (OL1) keeps AdamW moments task-only, preconditions the pressure
direction, projects its conflicting component relative to the adaptive task
direction, and caps the correction with a relative trust budget. This makes
the handling of pressure an explicit intervention. It does not guarantee
unchanged task-loss convergence or remove the need to specify pressure
strength and the trust budget.

\paragraph{From zeros to execution.}
TEAL's decoding kernels and the TwELL format and fused kernels of
\citet{cetin2026sparser} demonstrate that activation sparsity can yield
wall-clock gains in specified workloads~\citep{liu2024teal}. Their execution
designs make packing, memory access, and fusion part of the sparsity problem.
KernelBench evaluates correctness and speed in generated GPU
kernels~\citep{ouyang2025kernelbench}, while recent agentic kernel optimization
also studies irregular workloads such as mixture-of-experts and sparse
attention~\citep{luo2026agentic}. These results motivate kernel adaptation as
a separate test of whether a learned zero pattern is useful on a target
device. They do not supply a universal conversion from
$R_{\mathrm{model}}$ to speedup. For this study, logical sparsity and measured
execution performance therefore answer complementary questions.
